\bta{2017}


\begin{enumerate}
	%\renewcommand{\labelenumi}{\arabic{enumi}.}
	% A(\Alph) a(\alph) I(\Roman) i(\roman) 1(\arabic)
	%设定全局标号series=example	%引用全局变量resume=example
	%[topsep=-0.3em,parsep=-0.3em,itemsep=-0.3em,partopsep=-0.3em]
	%可使用leftmargin调整列表环境左边的空白长度 [leftmargin=0em]
	\item
	选择题 (本题满分 50 分，每小题 5 分)
\begin{enumerate}
	%\renewcommand{\labelenumi}{\arabic{enumi}.}
	% A(\Alph) a(\alph) I(\Roman) i(\roman) 1(\arabic)
	%设定全局标号series=example	%引用全局变量resume=example
	%[topsep=-0.3em,parsep=-0.3em,itemsep=-0.3em,partopsep=-0.3em]
	%可使用leftmargin调整列表环境左边的空白长度 [leftmargin=0em]
	\item
	函数 $f(x)=|x| \sin x^{2}$, 正确的结论是 \xzanswer{C}





\fourchoices
{在 $(-\infty,+\infty)$ 内有界}
{当 $x \rightarrow \infty$ 时, $f(x)$ 为无穷大}
{在 $(-\infty,+\infty)$ 内处处可导}
{当 $x \rightarrow \infty$ 时, $f(x)$ 极限存在}


\item
$ \lim\limits _{n \rightarrow \infty}\left(\frac{1}{n^{4}+n^{3}+n^{2}+1}+\frac{2}{n^{4}+n^{3}+n^{2}+2}+\cdots+\frac{n^{2}}{n^{4}+n^{3}+n^{2}+n^{2}}\right)= $ \xzanswer{B}


\fourchoices
{$ 0 $}
{$\frac{1}{2}$}
{$ 1 $}
{$\infty$}

\item
设 $f(x)=(x-1)(x-2)^{2}(x-3)^{3}$, 则 $f^{\prime}(1)+f^{\prime \prime}(2)+f^{\prime \prime \prime}(3)=$ \xzanswer{\text{无}}


\fourchoices
{$ 0 $}
{$-2$}
{$-8$}
{$-10$}

\item
设 $f(x)=x \sin x+\cos x-\frac{\pi}{2}$, 则下列命题中正确的是 \xzanswer{B}


\fourchoices
{$f(0)$ 是极大值, $f\left(\frac{\pi}{2}\right)$ 是极小值}
{$f(0)$ 是极小值, $f\left(\frac{\pi}{2}\right)$ 是极大值}
{$f(0)$ 是极大值, $f\left(\frac{\pi}{2}\right)$ 是极大值}
{$f(0)$ 是极小值, $f\left(\frac{\pi}{2}\right)$ 是极小值}


\item
求极限 $\lim\limits _{x \rightarrow 0}(1-3 x)^{\frac{\cos x^{2}}{\sin x}}=$ \xzanswer{C}


\fourchoices
{$e^{-1}$}
{$e^{-2}$}
{$e^{-3}$}
{$ 1 $}

\item
定积分 $\int_{0}^{\pi}(x \sin x)^{2} d x=$ \xzanswer{B}


\fourchoices
{$\frac{\pi^{2}}{3}-\frac{\pi}{4}$}
{$\frac{\pi^{3}}{6}-\frac{\pi}{4}$}
{$\frac{\pi^{3}}{6}-\frac{\pi^{2}}{4}$}
{$\frac{\pi^{3}}{6}+\frac{\pi}{4}$}

\item
设 $D_{K}$ 是圆域 $D=\left\{(x, y) \mid x^{2}+y^{2} \leqslant 1\right\}$ 在第 $k$ 象限的部分, $k=1,2,3,4$ 。 记 $I_{k}=\iint_{D_{k}}(y-x) d x d y$, 则正确的是 \xzanswer{D}


\fourchoices
{$I_{1}<0$}
{$I_{2}<0$}
{$I_{3}<0$}
{$I_{4}<0$}

\item
已知 $y_{1}(x)$ 和 $y_{2}(x)$ 是微分方程 $y^{\prime}+p(x) y=0$ 的两个不同的特解，则方程的通 解一定为 \xzanswer{D}


\fourchoices
{$y=C y_{1}(x)$}
{$y=Cy_{2}(x)$}
{$y=C\left(y_{1}(x)+y_{2}(x)\right)$}
{$y=C\left(y_{1}(x)-y_{2}(x)\right)$}

\item
设级数 $\sum\limits_{n=1}^{+\infty} u_{n}$ 收敛，则必收敛的级数为 \xzanswer{D}


\fourchoices
{$\sum\limits_{n=1}^{+\infty}(-1)^{n} \frac{u_{n}}{n}$}
{$\sum\limits_{n=1}^{+\infty} u_{n}^{2}$}
{$\sum\limits_{n=1}^{+\infty}\left(u_{2 n-1}-u_{2 n}\right)$}
{$\sum\limits_{n=1}^{+\infty}\left(u_{n}+u_{n+1}\right)$}

\item
设 $f^{\prime}\left(\sin ^{2} x\right)=\cos 2 x$, 且 $f(0)=1$, 则 $\int_{0}^{1} f(x) d x=$ \xzanswer{C}


\fourchoices
{$ 1 $}
{$\frac{1}{6}$}
{$\frac{7}{6}$}
{$\frac{1}{2}$}



\end{enumerate}


\banswer{
\begin{enumerate}
	%\renewcommand{\labelenumi}{\arabic{enumi}.}
	% A(\Alph) a(\alph) I(\Roman) i(\roman) 1(\arabic)
	\item
	\textbf{C}。取 $x_k=\sqrt{2k\pi+\frac{\pi}{2}}$，则 $f\left(x_k\right)=x_k\to+\infty$，故 $f$ 无界，A 错；
	再取 $t_k=\sqrt{k\pi}$，则 $f\left(t_k\right)=0$，故 $x\to\infty$ 时 $f$ 既不是无穷大，
	极限也不存在，B、D 均错。$x\neq0$ 时 $f$ 显然可导；在 $x=0$ 处
	\[
	\lim_{x\to0}\frac{|x|\sin x^2-0}{x}=\lim_{x\to0}\operatorname{sgn}(x)\sin x^2=0 ,
	\]
	故 $f'(0)=0$ 存在，$f$ 在 $(-\infty,+\infty)$ 内处处可导。

	\item
	\textbf{B}。记 $S_n=\sum\limits_{k=1}^{n^2}\frac{k}{n^4+n^3+n^2+k}$。
	由 $1\leqslant k\leqslant n^2$ 知分母介于 $n^4+n^3+n^2+1$ 与 $n^4+n^3+2n^2$ 之间，
	而 $\sum\limits_{k=1}^{n^2}k=\frac{n^2\left(n^2+1\right)}{2}$，故
	\[
	\frac{n^2\left(n^2+1\right)}{2\left(n^4+n^3+2n^2\right)}\leqslant S_n
	\leqslant\frac{n^2\left(n^2+1\right)}{2\left(n^4+n^3+n^2+1\right)} .
	\]
	两端均趋于 $\frac12$，由夹逼准则得 $S_n\to\frac12$。

	\item
	四个选项均不正确，正确值为 $2$。把 $f$ 在各根处写成“单因式乘余项”的形式：

	$f=(x-1)g$，$g=(x-2)^2(x-3)^3$，则 $f'(1)=g(1)=1\times(-8)=-8$；

	$f=(x-2)^2h$，$h=(x-1)(x-3)^3$，则 $f''(2)=2!\,h(2)=2\times(-1)=-2$；

	$f=(x-3)^3k$，$k=(x-1)(x-2)^2$，则 $f'''(3)=3!\,k(3)=6\times2=12$。

	故 $f'(1)+f''(2)+f'''(3)=-8-2+12=2$。

	\par\vspace{1mm}\noindent
	\textbf{注：}{\kaishu
	原答案作 D（$-10$），有误：$-10=-8+(-2)+0$，相当于把 $f'''(3)$ 误算成 $0$。
	$x=3$ 固然是三重根，但那只给出 $f(3)=f'(3)=f''(3)=0$，
	三阶导数反而是非零的 $3!\,k(3)=12$。
	正确值 $2$ 不在四个选项之中，本题选项有缺陷。
	}\par

	\item
	\textbf{B}。$f'(x)=\sin x+x\cos x-\sin x=x\cos x$，
	故 $f'(0)=f'\left(\frac{\pi}{2}\right)=0$。
	在 $x=0$ 附近 $\cos x>0$，$f'$ 由负变正，$f(0)$ 是极小值；
	在 $x=\frac{\pi}{2}$ 附近 $x>0$，$\cos x$ 由正变负，$f'$ 由正变负，
	$f\left(\frac{\pi}{2}\right)$ 是极大值。

	\item
	\textbf{C}。取对数，
	\[
	\ln L=\lim_{x\to0}\frac{\cos x^2}{\sin x}\ln(1-3x)
	=\lim_{x\to0}\frac{1\cdot(-3x)}{x}=-3 ,
	\]
	其中用了 $\ln(1+u)\sim u$、$\sin x\sim x$、$\cos x^2\to1$。故 $L=e^{-3}$。

	\item
	\textbf{B}。由 $\sin^2x=\frac{1-\cos2x}{2}$ 得
	\[
	\int_0^{\pi}x^2\sin^2x\,dx=\frac12\int_0^{\pi}x^2dx-\frac12\int_0^{\pi}x^2\cos2x\,dx .
	\]
	第一项为 $\frac{\pi^3}{6}$；第二项两次分部积分得
	$\int_0^{\pi}x^2\cos2x\,dx=\left[\frac{x^2\sin2x}{2}+\frac{x\cos2x}{2}-\frac{\sin2x}{4}\right]_0^{\pi}=\frac{\pi}{2}$，
	故原式 $=\frac{\pi^3}{6}-\frac{\pi}{4}$。

	\item
	\textbf{D}。$D_1$、$D_3$ 都关于直线 $y=x$ 对称，而 $y-x$ 在此对称下变号，
	故 $I_1=I_3=0$，A、C 错。在 $D_2$ 上 $x<0<y$，$y-x>0$，$I_2>0$，B 错。
	在 $D_4$ 上 $y<0<x$，$y-x<0$，故 $I_4<0$。

	\item
	\textbf{D}。方程的解集是一维线性空间：$y=Ce^{-\int p(x)dx}$，
	故 $y_1=C_1e^{-\int p}$，$y_2=C_2e^{-\int p}$，且 $C_1\neq C_2$。
	A、B 中若恰好 $C_1=0$（或 $C_2=0$），则 $Cy_1\equiv0$ 不是通解；
	C 中若 $C_2=-C_1$，则 $y_1+y_2\equiv0$ 也不行。
	只有 $y_1-y_2=\left(C_1-C_2\right)e^{-\int p}$ 因 $C_1\neq C_2$ 必非零解，
	故 $C\left(y_1-y_2\right)$ 一定是通解。

	\item
	\textbf{D}。设 $S_n=\sum\limits_{k=1}^{n}u_k\to S$，则
	\[
	\sum_{k=1}^{n}\left(u_k+u_{k+1}\right)=S_n+\left(S_{n+1}-u_1\right)\to2S-u_1 ,
	\]
	故 D 收敛。反例说明其余三个不必收敛：
	取 $u_n=\frac{(-1)^n}{\ln n}\ (n\geqslant2)$，$\sum u_n$ 由莱布尼茨判别法收敛，
	而 $\sum(-1)^n\frac{u_n}{n}=\sum\frac{1}{n\ln n}$ 发散，A 错；
	取 $u_n=\frac{(-1)^n}{\sqrt n}$，则 $\sum u_n^2=\sum\frac1n$ 发散，B 错；
	取 $u_n=\frac{(-1)^n}{n}$，则 $u_{2n-1}-u_{2n}=-\frac{1}{2n-1}-\frac{1}{2n}$，
	其绝对值与 $\frac1n$ 同阶，级数发散，C 错。

	\item
	\textbf{C}。令 $t=\sin^2x$，则 $\cos2x=1-2\sin^2x=1-2t$，
	且 $t$ 取遍 $[0,1]$，故 $f'(t)=1-2t$，积分得 $f(t)=t-t^2+C$；
	由 $f(0)=1$ 得 $C=1$，即 $f(x)=x-x^2+1$。于是
	\[
	\int_0^1f(x)dx=\frac12-\frac13+1=\frac76 .
	\]
\end{enumerate}
}


\item
(本题满分 10 分)
计算 $\lim\limits _{x \rightarrow 0} \frac{[\sin x-\sin (\sin x)] \sin x}{x^{4} \cos x^{2}}$。

\banswer{
	由 $u-\sin u=\frac{u^3}{6}+o\left(u^3\right)$，取 $u=\sin x$，并注意 $\sin x\sim x$，得
	\[
	\sin x-\sin(\sin x)=\frac{\sin^3x}{6}+o\left(x^3\right)\sim\frac{x^3}{6} .
	\]
	又 $\sin x\sim x$，$\cos x^2\to1$，故
	\[
	\lim_{x\to0}\frac{\left[\sin x-\sin(\sin x)\right]\sin x}{x^4\cos x^2}
	=\lim_{x\to0}\frac{\frac{x^3}{6}\cdot x}{x^4}=\frac16 .
	\]

	\par\vspace{1mm}\noindent
	\textbf{注：}{\kaishu
	这里用的是“等价无穷小可以整块替换乘积因子”的规则：
	$\sin x-\sin(\sin x)$ 与 $\sin x$ 都是被乘的因子，故可分别替换；
	若把它们拆成 $\sin x$ 与 $\sin(\sin x)$ 之差再各自替换成 $x$，就会得到错误的 $0$。
	}\par
}



\item
(本题满分 10 分)
证明: $\lim\limits _{n \rightarrow \infty} \int_{0}^{\pi} \sin ^{n} t d t=0$。


\banswer{
	由 $\sin(\pi-t)=\sin t$ 知
	$\int_0^{\pi}\sin^nt\,dt=2\int_0^{\pi/2}\sin^nt\,dt=2I_n$，
	故只需证 $I_n=\int_0^{\pi/2}\sin^nt\,dt\to0$。

	\textbf{解法一（分割区间）.}
	任给 $\varepsilon>0$，先取 $\delta\in\left(0,\frac{\pi}{2}\right)$ 使 $\delta<\frac{\varepsilon}{4}$，
	再把积分区间分成两段：
	\begin{align*}
	\int_0^{\pi/2-\delta}\sin^nt\,dt&\leqslant\frac{\pi}{2}\left(\cos\delta\right)^n,\\
	\int_{\pi/2-\delta}^{\pi/2}\sin^nt\,dt&\leqslant\delta<\frac{\varepsilon}{4} .
	\end{align*}
	第一式用了：$t\in\left[0,\frac{\pi}{2}-\delta\right]$ 时
	$0\leqslant\sin t\leqslant\sin\left(\frac{\pi}{2}-\delta\right)=\cos\delta<1$。
	由 $0<\cos\delta<1$ 知 $\left(\cos\delta\right)^n\to0$，故存在 $N$，
	当 $n>N$ 时第一式 $<\frac{\varepsilon}{4}$。于是 $n>N$ 时
	$0<2I_n<\varepsilon$，即 $\int_0^{\pi}\sin^nt\,dt\to0$。

	\textbf{解法二（Wallis 递推）.}
	分部积分给出 $I_n=\frac{n-1}{n}I_{n-2}\ (n\geqslant2)$，即 $nI_n=(n-1)I_{n-2}$。
	两端乘 $I_{n-1}$ 得 $nI_nI_{n-1}=(n-1)I_{n-1}I_{n-2}$，
	故 $nI_nI_{n-1}$ 与 $n$ 无关，等于 $1\cdot I_1I_0=1\times\frac{\pi}{2}=\frac{\pi}{2}$。
	又 $0\leqslant\sin t\leqslant1$ 使 $\{I_n\}$ 单调不增且为正，故 $I_n\leqslant I_{n-1}$，
	\[
	I_n^2\leqslant I_nI_{n-1}=\frac{\pi}{2n}
	\ \Longrightarrow\ 0<I_n\leqslant\sqrt{\frac{\pi}{2n}}\to0 .
	\]
	从而 $\int_0^{\pi}\sin^nt\,dt=2I_n\to0$（且衰减速度为 $\sqrt{2\pi/n}$）。
}


\item
(本题满分 10 分)
把函数 $f(x)=(x-1)^{2}$ 在 $(0,1)$ 上展成余弦级数，并求 $\sum\limits_{n=1}^{+\infty} \frac{1}{n^{2}}$。


\banswer{
	作偶延拓再作周期为 $2$ 的延拓，则 $l=1$，只有余弦项：
	\[
	a_n=\frac{2}{l}\int_0^{l}f(x)\cos\frac{n\pi x}{l}dx=2\int_0^1(x-1)^2\cos n\pi x\,dx .
	\]

	\textbf{定 $a_0$.} $a_0=2\int_0^1(x-1)^2dx=2\times\frac13=\frac23$，故常数项 $\frac{a_0}{2}=\frac13$。

	\textbf{定 $a_n\ (n\geqslant1)$.} 令 $u=x-1$，注意
	$\cos\left[n\pi(u+1)\right]=(-1)^n\cos n\pi u$，且 $u^2\cos n\pi u$ 为偶函数，
	\[
	a_n=2(-1)^n\int_{-1}^{0}u^2\cos n\pi u\,du=2(-1)^n\int_0^1u^2\cos n\pi u\,du .
	\]
	两次分部积分得 $\int_0^1u^2\cos n\pi u\,du=\frac{2(-1)^n}{n^2\pi^2}$，故
	$a_n=\frac{4}{n^2\pi^2}$。于是在 $(0,1)$ 上
	\[
	f(x)=\frac13+\frac{4}{\pi^2}\sum_{n=1}^{\infty}\frac{\cos n\pi x}{n^2} .
	\]

	\textbf{求和.} 偶延拓后的函数在 $x=0$ 处连续，级数在该点收敛到 $f(0)=1$，故
	\[
	1=\frac13+\frac{4}{\pi^2}\sum_{n=1}^{\infty}\frac{1}{n^2}
	\ \Longrightarrow\
	\sum_{n=1}^{\infty}\frac{1}{n^2}=\frac{\pi^2}{4}\times\frac23=\frac{\pi^2}{6} .
	\]
}


\item
(本题满分 10 分)
计算曲面积分
\[
I=\iint_{\Sigma} 2 x^{3} d y d z+2 y^{3} d z d x+3\left(z^{2}-2\right) d x d y
\]
其中 $\Sigma$ 为曲面 $z=1-x^{2}-y^{2}(z \geqslant 0)$ 的上侧。


\banswer{
	\textbf{解法一（高斯公式）.} 补上底面
	$\Sigma_1:\ z=0,\ x^2+y^2\leqslant1$ 并取下侧，
	则 $\Sigma+\Sigma_1$ 是旋转抛物体 $\Omega$ 的外侧闭曲面。
	$\operatorname{div}\bm{F}=6x^2+6y^2+6z$，用柱坐标（$0\leqslant r\leqslant1$，$0\leqslant z\leqslant1-r^2$）：
	\begin{align*}
	\iiint_{\Omega}6r^2\,dV&=2\pi\int_0^1 6r^3\left(1-r^2\right)dr=\pi,\\
	\iiint_{\Omega}6z\,dV&=2\pi\int_0^1 3r\left(1-r^2\right)^2dr=\pi,
	\end{align*}
	故 $\oiint_{\Sigma+\Sigma_1}=2\pi$。在 $\Sigma_1$ 上 $dydz=dzdx=0$，$z=0$，取下侧，
	\[
	\iint_{\Sigma_1}=-\iint_{D}3(0-2)\,dxdy=6\pi .
	\]
	于是 $I=2\pi-6\pi=-4\pi$。

	\textbf{解法二（直接投影）.} 取上侧时
	$d\bm{S}=\left(-z_x,-z_y,1\right)dxdy=(2x,2y,1)dxdy$，
	在 $D:\ x^2+y^2\leqslant1$ 上
	\[
	I=\iint_{D}\left[4x^4+4y^4+3\left(\left(1-r^2\right)^2-2\right)\right]dxdy ,
	\]
	其中 $r^2=x^2+y^2$。用极坐标并注意
	$\int_0^{2\pi}\cos^4\theta\,d\theta=\int_0^{2\pi}\sin^4\theta\,d\theta=\frac{3\pi}{4}$，
	得 $\iint_D4\left(x^4+y^4\right)dxdy=4\times\frac{3\pi}{2}\times\frac16=\pi$，
	而 $3\iint_D\left[\left(1-r^2\right)^2-2\right]dxdy=3\times2\pi\times\left(\frac16-1\right)=-5\pi$，
	合计 $I=\pi-5\pi=-4\pi$，与解法一一致。

	\par\vspace{1mm}\noindent
	\textbf{注：}{\kaishu
	数值积分给出 $I\approx-12.5664=-4\pi$，与上述结果吻合。
	}\par
}



\item
(本题满分 10 分)
计算 $\iint_{\Omega} \frac{(x+y) \ln \left(1+\frac{y}{x}\right)}{\sqrt{1-x-y}} d x d y$, 这里 $\Omega$ 是由直线 $x+y=1, x=0, y=0$ 所围成的 三角形区域。

\banswer{
	被积函数只依赖 $x+y$ 与 $\frac{y}{x}$，故作变换
	\[
	\left\{\begin{array}{l}t=x+y\\ s=x\end{array}\right.
	\quad\Longrightarrow\quad x=s,\ y=t-s,\ \left|\frac{\partial(x,y)}{\partial(t,s)}\right|=1 .
	\]
	三角形 $\Omega$ 变成 $0\leqslant t\leqslant1$，$0\leqslant s\leqslant t$，且
	$\ln\left(1+\frac{y}{x}\right)=\ln\frac{x+y}{x}=\ln\frac{t}{s}$，于是
	\[
	I=\int_0^1\frac{t}{\sqrt{1-t}}\left(\int_0^t\ln\frac{t}{s}\,ds\right)dt .
	\]

	\textbf{内层积分.} $\int_0^t\left(\ln t-\ln s\right)ds=t\ln t-\left[s\ln s-s\right]_0^t=t$。

	\textbf{外层积分.} 代入并令 $u=1-t$：
	\[
	I=\int_0^1\frac{t^2}{\sqrt{1-t}}dt=\int_0^1\frac{(1-u)^2}{\sqrt u}du
	=2-\frac43+\frac25=\frac{16}{15} .
	\]

	\par\vspace{1mm}\noindent
	\textbf{注：}{\kaishu
	$s\to0^{+}$ 时 $\ln\frac{t}{s}\to+\infty$、$t\to1^{-}$ 时 $\frac{1}{\sqrt{1-t}}\to+\infty$，
	都是可积的瑕点，上述计算合法。数值积分给出 $I\approx1.06667=\frac{16}{15}$。
	}\par
}


\newpage
\item
(本题满分 10 分)
求微分方程 $y^{\prime \prime}\left(x+\left(y^{\prime}\right)^{2} \right)=y^{\prime}$ 满足初始条件 $y(1)=y^{\prime}(1)=1$ 的特解。


\banswer{
	方程不显含 $y$，令 $p=y'$。由初值 $p(1)=1>0$ 知在 $x=1$ 附近 $p\neq0$，
	故可把 $x$ 看成 $p$ 的函数：$y''=\frac{dp}{dx}=\left(\frac{dx}{dp}\right)^{-1}$，代入原方程得
	\[
	\frac{x+p^2}{\frac{dx}{dp}}=p
	\quad\Longrightarrow\quad
	\frac{dx}{dp}-\frac{x}{p}=p ,
	\]
	这是以 $x$ 为未知函数的一阶线性方程。乘积分因子 $\frac1p$：
	\[
	\frac{d}{dp}\left(\frac{x}{p}\right)=1
	\quad\Longrightarrow\quad
	\frac{x}{p}=p+C_1,\qquad x=p^2+C_1p .
	\]

	\textbf{定 $C_1$.} 由 $x=1$ 时 $p=1$ 得 $1=1+C_1$，即 $C_1=0$，故 $x=p^2$，
	取正根 $p=y'=\sqrt{x}$。

	\textbf{定 $C_2$.} 积分得 $y=\frac23x^{3/2}+C_2$，由 $y(1)=1$ 得 $C_2=\frac13$。所求特解为
	\[
	y=\frac23x^{3/2}+\frac13 .
	\]
	（验证：$y'=\sqrt x$，$y''=\frac{1}{2\sqrt x}$，
	$y''\left(x+y'^2\right)=\frac{2x}{2\sqrt x}=\sqrt x=y'$。）
}


\item
(本题满分 10 分)
计算曲线积分 $I=\oint_{L} \frac{x d y-y d x}{4 x^{2}+y^{2}}$, 其中 $L$ 是以点 $(1,0)$ 为中心, $ 2  $为半径的圆周, 取 逆时针方向。


\banswer{
	记 $P=\frac{-y}{4x^2+y^2}$，$Q=\frac{x}{4x^2+y^2}$。除原点外
	\[
	\frac{\partial Q}{\partial x}=\frac{y^2-4x^2}{\left(4x^2+y^2\right)^2}
	=\frac{\partial P}{\partial y} .
	\]
	而原点到圆心 $(1,0)$ 的距离为 $1<2$，故奇点 $(0,0)$ 在 $L$ 内部，不能直接用格林公式。

	\textbf{挖去奇点.} 取充分小的 $\varepsilon>0$，作椭圆
	$l_\varepsilon:\ 4x^2+y^2=\varepsilon^2$（含于 $L$ 内），取逆时针方向。
	在 $L$ 与 $l_\varepsilon$ 所夹的复连通区域上 $Q_x=P_y$，由格林公式得 $I=\oint_{l_\varepsilon}$。

	\textbf{算小椭圆.} 取参数 $x=\frac{\varepsilon}{2}\cos\theta$，$y=\varepsilon\sin\theta$，$\theta:0\to2\pi$，则
	\[
	x\,dy-y\,dx=\frac{\varepsilon^2}{2}\left(\cos^2\theta+\sin^2\theta\right)d\theta
	=\frac{\varepsilon^2}{2}d\theta ,
	\]
	而分母恒为 $\varepsilon^2$，故
	\[
	I=\int_0^{2\pi}\frac12\,d\theta=\pi .
	\]

	\par\vspace{1mm}\noindent
	\textbf{注：}{\kaishu
	必须换成与分母配套的椭圆（而不是圆），才能使分母变成常数。
	沿原圆周直接数值积分也得 $I\approx3.14159=\pi$。
	}\par
}


\item
(本题满分 10 分)
设 $x \geqslant 0, a>0$, 若 $\sqrt{x+a}-\sqrt{x}=\frac{a}{2 \sqrt{x+\phi_{a}(x)}}$ 成立，证明: $\frac{a}{4} \leqslant \phi_{a}(x) \leqslant \frac{a}{2}$ 且 $\lim\limits _{x \rightarrow 0} \phi_{a}(x)=\frac{a}{4}$ 和 $\lim\limits _{x \rightarrow \infty} \phi_{a}(x)=\frac{a}{2}$。

\banswer{
	\textbf{解出 $\phi_a$.} 由 $\sqrt{x+a}-\sqrt x=\frac{a}{\sqrt{x+a}+\sqrt x}$ 及题设知
	$2\sqrt{x+\phi_a(x)}=\sqrt{x+a}+\sqrt{x}$。两端平方：
	\[
	4\left(x+\phi_a\right)=2x+a+2\sqrt{x^2+ax} ,
	\]
	故
	\[
	\phi_a(x)=\frac{a}{4}+\frac{1}{2}\left(\sqrt{x^2+ax}-x\right)
	=\frac{a}{4}+\frac{1}{2}h(x) .
	\]

	\textbf{估计 $h$.} 分子有理化：对 $x>0$
	\[
	h(x)=\frac{ax}{\sqrt{x^2+ax}+x}=\frac{a}{\sqrt{1+\frac{a}{x}}+1} ,
	\]
	且 $h(0)=0$。因 $\frac{a}{x}$ 随 $x$ 严格递减，$h$ 在 $(0,+\infty)$ 上严格递增，
	并有 $0\leqslant h(x)<\frac{a}{2}$。

	\textbf{结论.} 代回即得
	\[
	\frac{a}{4}\leqslant\phi_a(x)<\frac{a}{4}+\frac{a}{4}=\frac{a}{2} ,
	\]
	当然更有 $\frac a4\leqslant\phi_a(x)\leqslant\frac a2$。又由上式
	\[
	\lim_{x\to0^{+}}\phi_a(x)=\frac a4+0=\frac a4,\qquad
	\lim_{x\to+\infty}\phi_a(x)=\frac a4+\frac a4=\frac a2 .
	\]

	\par\vspace{1mm}\noindent
	\textbf{注：}{\kaishu
	$\phi_a$ 恰是拉格朗日中值定理
	$\sqrt{x+a}-\sqrt{x}=\frac{a}{2\sqrt{\xi}}$ 中的中值 $\xi=x+\phi_a(x)$ 的位置；
	结论说明：区间很短（$x\to\infty$）时中值趋于区间中点，
	而 $x\to0$ 时因 $\sqrt{\ }$ 在端点附近变化剧烈，中值被推向左端的 $\frac14$ 处。
	上界 $\frac a2$ 其实取不到。
	}\par
}


\item
(本题满分 10 分)
函数 $f(x)$ 在 $[0,2]$ 上连续, 且为增函数，证明:
\[
\int_{0}^{2} f(x) d x \leq \int_{0}^{2} x f(x) d x
\]

\banswer{
	只需证 $J=\int_0^2(x-1)f(x)dx\geqslant0$。注意 $\int_0^2(x-1)dx=0$，故可以减去常数 $f(1)$：
	\[
	J=\int_0^2(x-1)\left[f(x)-f(1)\right]dx .
	\]
	因 $f$ 单调增：$x>1$ 时 $x-1>0$ 且 $f(x)-f(1)\geqslant0$；
	$x<1$ 时 $x-1<0$ 且 $f(x)-f(1)\leqslant0$。两种情形下被积函数都 $\geqslant0$，
	故 $J\geqslant0$，即
	\[
	\int_0^2f(x)dx\leqslant\int_0^2xf(x)dx .
	\]

	\par\vspace{1mm}\noindent
	\textbf{注：}{\kaishu
	“减去 $f(1)$ ”这一步是关键：它把符号不定的 $(x-1)f(x)$ 变成了处处非负的
	$(x-1)\left[f(x)-f(1)\right]$。此法对任意区间 $[a,b]$ 与权 $x-\frac{a+b}{2}$ 都适用。
	}\par
}

\item
(本题满分 10 分)
设 $f(x)$ 在 $[a, b]$ 上可导且 $f^{\prime}(x)$ 连续, $f(a)=0$, 这里 $a<b$。 证明:
\[
\int_{a}^{b} \left[f(x)\right] ^{2} d x \leq \frac{(b-a)^{2}}{2} \int_{a}^{b}\left[f^{\prime}(x)\right]^{2} d x
\]

\banswer{
	由 $f(a)=0$ 及牛顿--莱布尼茨公式，对每个 $x\in[a,b]$ 有
	$f(x)=\int_a^xf'(t)dt$。用柯西--施瓦茨不等式（积分形式）：
	\[
	f^2(x)=\left(\int_a^x1\cdot f'(t)dt\right)^2
	\leqslant(x-a)\int_a^x\left[f'(t)\right]^2dt .
	\]
	再把右端的积分上限放大到 $b$（被积函数非负）：
	\[
	f^2(x)\leqslant(x-a)\int_a^b\left[f'(t)\right]^2dt .
	\]
	两端对 $x$ 在 $[a,b]$ 上积分，并注意 $\int_a^b(x-a)dx=\frac{(b-a)^2}{2}$，得
	\[
	\int_a^b f^2(x)dx\leqslant\frac{(b-a)^2}{2}\int_a^b\left[f'(x)\right]^2dx .
	\]

	\par\vspace{1mm}\noindent
	\textbf{注：}{\kaishu
	这是 Poincaré 型不等式。常数 $\frac{(b-a)^2}{2}$ 并非最优
	（最优常数是 $\frac{4(b-a)^2}{\pi^2}$，由 $f(x)=\sin\frac{\pi(x-a)}{2(b-a)}$ 达到），
	但本题只要求证明上式。
	}\par
}

\end{enumerate}

